\subsection{DIRECT LIMITS OF MODULES}

\BourbakiRunInHeading{Henceforth I is assumed to be right directed.}

Let \((\mathbf{A}_{\alpha}, \phi_{\beta\alpha})\) be a direct system of rings (I, § 10, no. 3), \((\mathbf{E}_{\alpha}, f_{\beta\alpha})\) a direct system of commutative groups (written additively) (I, § 10, no. 3) and suppose that each \(E_{\alpha}\) has a left \(A_{\alpha}\) -module structure; further, suppose that, for \(\alpha \leqslant \beta\) , \((f_{\beta\alpha}, \phi_{\beta\alpha})\) is a dimorphism of \(E_{\alpha}\) into \(E_{\beta}\) (§ 1, no. 13), in other words that

\[
f _ {\beta \alpha} (\lambda_ {\alpha} x _ {\alpha}) = \phi_ {\beta \alpha} (\lambda_ {\alpha}) f _ {\beta \alpha} (x _ {\alpha})\tag{4}
\]

for \(x_{\alpha} \in \mathrm{E}_{\alpha}, \lambda_{\alpha} \in \mathrm{A}_{\alpha}\) ; then \(\mathbf{E} = \varinjlim \mathbf{E}_{\alpha}\) has a left module structure over \(\mathbf{A} = \varinjlim \mathbf{A}_{\alpha}\) (I, § 10, no. 4). For all \(\alpha \in \mathbf{I}\) , let \(f_{\alpha}: \mathrm{E}_{\alpha} \to \mathrm{E}\) , \(\phi_{\alpha}: \mathrm{A}_{\alpha} \to \mathrm{A}\) be the canonical mappings; then \((f_{\alpha}, \phi_{\alpha})\) is a dimorphism of \(\mathbf{E}_{\alpha}\) into \(\mathbf{E}\) . We shall say that \((\mathbf{E}_{\alpha}, f_{\beta \alpha})\) is a direct system of left \(\mathbf{A}_{\alpha}\) -modules and that the A-module \(\mathbf{E}\) is its direct limit.

Let \((\mathbf{E}_{\alpha}^{\prime}, f_{\beta \alpha}^{\prime})\) be another direct system of left \(A_{\alpha}\) -modules and, for all \(\alpha\) , let \(u_{\alpha}: \mathrm{E}_{\alpha}' \to \mathrm{E}_{\alpha}\) be an \(\mathbf{A}_{\alpha}\) -linear mapping, these mappings forming a direct system; then \(u = \varinjlim u_{\alpha}\) is an A-linear mapping of \(\varinjlim \mathrm{E}_{\alpha}'\) into \(\varinjlim \mathrm{E}_{\alpha}\) .

Moreover:

PROPOSITION 3. Let \((\mathbf{E}_{\alpha}, f_{\beta \alpha})\) , \((\mathbf{E}_{\alpha}', f_{\beta \alpha}')\) , \((\mathbf{E}_{\alpha}'', f_{\beta \alpha}'')\) be three direct systems of \(\mathbf{A}_{\alpha}\) -modules and \((u_{\alpha})\) , \((v_{\alpha})\) two direct systems of \(\mathbf{A}_{\alpha}\) -linear mappings such that the sequences

\[
\mathrm{E} _ {\alpha} ^ {\prime} \xrightarrow {u _ {\alpha}} \mathrm{E} _ {\alpha} \xrightarrow {v _ {\alpha}} \mathrm{E} _ {\alpha} ^ {\prime \prime}
\]

are exact for all \(\alpha\) . Then, writing \(u = \varinjlim u_{\alpha}\) , \(v = \varinjlim v_{\alpha}\) , the sequence

\[
\lim _ {\longrightarrow} \mathrm{E} _ {\alpha} ^ {\prime} \xrightarrow {u} \lim _ {\longrightarrow} \mathrm{E} _ {\alpha} \xrightarrow {v} \lim _ {\longrightarrow} \mathrm{E} _ {\alpha} ^ {\prime \prime}
\]

is exact.

\(u(\varinjlim \mathrm{E}_{\alpha}^{\prime}) = \varinjlim u_{\alpha}(\mathrm{E}_{\alpha}^{\prime})\) and \(v^{-1}(0) = \varinjlim v_{\alpha}^{-1}(0)\) (Set Theory, III, § 7, no. 6, Corollary to Proposition 7).

Loosely speaking, Proposition 3 can also be expressed by saying that passing to the direct limit preserves exactness.

PROPOSITION 4. Let \((\mathbf{E}_{\alpha}, f_{\beta \alpha})\) be a direct system of \(\mathbf{A}_{\alpha}\) -modules, \(\mathbf{E} = \varinjlim \mathbf{E}_{\alpha}\) its direct limit and \(\phi_{\alpha}: \mathbf{A}_{\alpha} \to \mathbf{A}\) and \(f_{\alpha}: \mathbf{E}_{\alpha} \to \mathbf{E}\) the canonical mappings for all \(\alpha \in \mathbf{I}\) . If, for all \(\alpha \in \mathbf{I}\) , \(\mathbf{S}_{\alpha}\) is a generating system of \(\mathbf{E}_{\alpha}\) , then \(\mathbf{S} = \bigcup_{\alpha \in \mathbf{I}} f_{\alpha}(\mathbf{S}_{\alpha})\) is a generating system of \(\mathbf{E}\) .

Every \(x \in \mathbf{E}\) is of the form \(f_{\alpha}(x_{\alpha})\) for some \(\alpha \in \mathbf{I}\) and some \(x_{\alpha} \in \mathbf{E}_{\alpha}\) and by hypothesis \(x = \sum_{i} \lambda_{\alpha}^{(i)} y_{\alpha}^{(i)}\) , where \(\lambda_{\alpha}^{(i)} \in \mathbf{A}_{\alpha}\) and \(y_{\alpha}^{(i)} \in \mathbf{S}_{\alpha}\) ; writing \(\lambda^{(i)} = \phi_{\alpha}(\lambda_{\alpha}^{(i)})\) , \(y^{(i)} = f_{\alpha}(y_{\alpha}^{(i)})\) , we obtain \(x = \sum_{i} \lambda^{(i)} y^{(i)}\) .

PROPOSITION 5. With the hypotheses and notation of Proposition 4, suppose that for all \(\alpha \in I\) , \(E_{\alpha}\) is the direct sum of a family \((M_{\alpha}^{\lambda})_{\alpha \in L}\) of submodules (the indexing set \(L\) being independent of \(\alpha\) ) and that \(f_{\beta \alpha}(M_{\alpha}^{\lambda}) \subset M_{\alpha}^{\lambda}\) for \(\alpha \leqslant \beta\) and for all \(\lambda \in L\) . Then \(E\) is the direct sum of the submodules \(M^{\lambda} = \varinjlim_{\alpha} M_{\alpha}^{\lambda} (\lambda \in L)\) .

It follows from Proposition 4 that E is the sum of the \(\mathbf{M}^{\lambda}\) . Let \((y^{\lambda})_{\lambda \in L}\) be a family such that \(y^{\lambda} \in \mathbf{M}^{\lambda}\) for all \(\lambda \in L\) and whose support is finite and suppose that \(\sum_{\lambda} y^{\lambda} = 0\) . By virtue of Set Theory, III, §7, no. 5, Lemma 1, there exist an \(\alpha \in I\) and a family \((x_{\alpha}^{\lambda})_{\alpha \in L}\) of finite support consisting of elements of \(\mathbf{E}_{\alpha}\) such that \(x_{\alpha}^{\lambda} \in \mathbf{M}_{\alpha}^{\lambda}\) and \(y^{\lambda} = f_{\alpha}(x_{\alpha}^{\lambda})\) for all \(\lambda \in L\) . The relation \(f_{\alpha}\left(\sum_{\lambda \in L} x_{\alpha}^{\lambda}\right) = 0\) implies the existence of a \(\beta \geqslant \alpha\) such that \(f_{\beta \alpha}\left(\sum_{\lambda \in L} x_{\alpha}^{\lambda}\right) = 0\) (Set Theory, III, § 7, no. 5, Lemma 1), which may be written as \(\sum_{\lambda \in L} x_{\beta}^{\lambda} = 0\) , where \(x_{\beta}^{\lambda} = f_{\beta \alpha}(x_{\alpha}^{\lambda}) \in \mathrm{M}_{\beta}^{\lambda}\) by hypothesis; hence \(x_{\beta}^{\lambda} = 0\) for all \(\lambda \in \mathbf{L}\) and therefore \(y^{\lambda} = f_{\beta}(x_{\beta}^{\lambda}) = 0\) for all \(\lambda \in \mathbf{L}\) , which proves that the sum of the \(\mathbf{M}^{\lambda}\) is direct.

COROLLARY. Let \((\mathbf{P}_{\alpha})\) be a direct system of subsets of \(\mathbf{E}_{\alpha}\) and let \(\mathbf{P} = \varinjlim \mathbf{P}_{\alpha}\) . If, for all \(\alpha \in \mathbf{I}\) , \(\mathbf{P}_{\alpha}\) is a free subset (resp. basis) of \(\mathbf{E}_{\alpha}\) , then \(\mathbf{P}\) is a free subset (resp. basis) of \(\mathbf{E}\) .

The second assertion follows immediately from the first and Proposition 4. It is therefore sufficient to prove that if the \(\mathbf{P}_{\alpha}\) are free every subset \(\{y^{(i)}\}_{1\leqslant i\leqslant n}\) consisting of distinct elements of \(\mathbf{P}\) , is free. There exists an \(\alpha \in \mathbf{I}\) and elements \(x_{\alpha}^{(i)}\in \mathbf{P}_{\alpha}\) such that \(y^{(i)} = f_{\alpha}(x_{\alpha}^{(i)})\) for \(1\leqslant i\leqslant n\) (Set Theory, III, §7, no. 5, Lemma 1); if \(\sum_{i}\lambda^{(i)}y^{(i)} = 0\) , it may be assumed that \(\lambda^{(i)} = \phi_{\alpha}(\lambda_{\alpha}^{(i)})\) for \(1\leqslant i\leqslant n\) and hence \(f_{\alpha}\left(\sum_{i}\lambda_{\beta}^{(i)}x_{\beta}^{(i)}\right) = 0\) ; this implies \(\sum_{i}\lambda_{\beta}^{(i)}x_{\beta}^{(i)} = 0\) for some \(\beta \geqslant \alpha\) , where \(\lambda_{\beta}^{(i)} = \phi_{\beta \alpha}(\lambda_{\alpha}^{(i)})\) , \(x_{\beta}^{(i)} = f_{\beta \alpha}(x_{\alpha}^{(i)})\) and the \(x_{\beta}^{(i)}\) belong to \(\mathbf{P}_{\beta}\) and are distinct since \(y^{(i)} = f_{\beta}(x_{\beta}^{(i)})\) ; then \(\lambda_{\beta}^{(i)} = 0\) for \(1\leqslant i\leqslant n\) , whence

\[
\lambda^ {(i)} = \phi_ {\beta} (\lambda_ {\beta} ^ {(i)}) = 0
\]

for \(1 \leqslant i \leqslant n\) .

Suppose now that all the rings \(\mathbf{A}_{\alpha}\) are equal to the same ring A and the \(\phi_{\beta \alpha}\) to \(1_{\mathbf{A}}\) ; then, for every direct system \((\mathbf{E}_{\alpha}, f_{\beta \alpha})\) of A-modules, \(\mathbf{E} = \varinjlim \mathbf{E}_{\alpha}\) is an A-module. Let F be an A-module and for all \(\alpha\) let \(u_{\alpha}: \mathbf{E}_{\alpha} \to \mathbf{F}\) be an A-linear mapping such that \((u_{\alpha})\) is a direct system of mappings; then \(u = \varinjlim u_{\alpha}\) is an A-linear mapping of E into F. Conversely, for every A-linear mapping \(v: \varinjlim \mathbf{E}_{\alpha} \to \mathbf{F}\) , the family of \(v_{\alpha} = v \circ f_{\alpha}\) is a direct system of A-linear mappings such that \(v = \varinjlim v_{\alpha}\) . On the other hand we note that for \(\alpha \leqslant \beta\) the mapping

\[
\operatorname{Hom} (f _ {\beta \alpha}, 1 _ {F}) = \bar {f} _ {\alpha \beta}: \operatorname{Hom} _ {A} (E _ {\beta}, F) \rightarrow \operatorname{Hom} _ {A} (E _ {\alpha}, F)
\]

is a Z-module homomorphism such that \((\mathrm{Hom}_{\mathbf{A}}(\mathrm{E}_{\alpha},\mathrm{F}),\bar{f}_{\alpha\beta})\) is an inverse system of Z-modules; as \(\bar{f}_{\alpha\beta}(v_{\beta}) = v_{\beta} \circ f_{\beta\alpha}\) , the above remarks can be expressed as follows:

PROPOSITION 6. For every direct system \((\mathbf{E}_{\alpha}, f_{\beta \alpha})\) of A-modules and every A-module F, the canonical mapping \(u \mapsto (u \circ f_{\alpha})\) is a Z-module isomorphism

\[
d _ {\mathrm{F}}: \operatorname{Hom} _ {\mathrm{A}} (\varinjlim \mathrm{E} _ {\alpha}, \mathrm{F}) \rightarrow \varprojlim \operatorname{Hom} _ {\mathrm{A}} (\mathrm{E} _ {\alpha}, \mathrm{F}).\tag{5}
\]

COROLLARY 1. For every A-module homomorphism \(v: \mathbf{F} \to \mathbf{F}'\) , the

\[
\bar {v} _ {\alpha} = \operatorname{Hom} \left(\mathrm{l} _ {\mathrm{E} _ {\alpha}}, v\right): \operatorname{Hom} \left(\mathrm{E} _ {\alpha}, \mathrm{F}\right)\rightarrow \operatorname{Hom} \left(\mathrm{E} _ {\alpha}, \mathrm{F} ^ {\prime}\right)
\]

form an inverse system of Z-linear mappings and the diagram

\[
\begin{array}{c} \operatorname{Hom} (\varinjlim \mathrm{E} _ {\alpha}, \mathrm{F}) \xrightarrow {d _ {\mathrm{F}}} \varprojlim \operatorname{Hom} (\mathrm{E} _ {\alpha}, \mathrm{F}) \\ \operatorname{Hom} (1 _ {\mathrm{E}}, v) \Bigg \downarrow \qquad \qquad \qquad \qquad \qquad \qquad \qquad \qquad \qquad \qquad \qquad \qquad \qquad \qquad \qquad \varprojlim_ {\overline {{v}} _ {\alpha}} \\ \operatorname{Hom} (\varinjlim \mathrm{E} _ {\alpha}, \mathrm{F} ^ {\prime}) \xrightarrow {d _ {\mathrm{F} ^ {\prime}}} \varprojlim \operatorname{Hom} (\mathrm{E} _ {\alpha}, \mathrm{F} ^ {\prime}) \end{array}\tag{6}
\]

is commutative.

For all \(u \in \operatorname{Hom}(\varinjlim E_{\alpha}, F)\) , \(d_{F'}(v \circ u) = (v \circ u \circ f_{\alpha})\) by definition and the commutativity of diagram (6) then follows immediately from the definitions.

COROLLARY 2. If \((\mathbf{E}_{\alpha}, f_{\beta \alpha})\) is a direct system of left A-modules and \(\mathbf{E} = \varinjlim \mathbf{E}_{\alpha}\) , \((\mathbf{E}_{\alpha}^{*}, {}^t f_{\beta \alpha})\) is an inverse system of right A-modules and \(\varprojlim \mathbf{E}_{\alpha}^{*}\) is canonically isomorphic to \(\mathbf{E}^{*}\) .

Remark. Let E be an A-module and \((\mathrm{M}_{\alpha})_{\alpha\in I}\) an increasing family of submodules of E such that E is the union of the \(M_{\alpha}\) ; if \(j_{\beta\alpha}:M_{\alpha}\to M_{\beta}\) (for \(\alpha\leqslant\beta\) ) and \(j_{\alpha}:M_{\alpha}\to E\) are the canonical injections, it is immediate that \(j=\varinjlim j_{\alpha}\) is an isomorphism of \(\varinjlim M_{\alpha}\) onto E (Set Theory, III, §7, no. 6, Remark 1). In particular, every A-module is the direct limit of the right directed family of its finitely generated submodules.

